
\documentclass{article}
\usepackage{colortbl}
\usepackage{makecell}
\usepackage{multirow}
\usepackage{supertabular}

\begin{document}

\newcounter{utterance}

\twocolumn

{ \footnotesize  \setcounter{utterance}{1}
\setlength{\tabcolsep}{0pt}
\begin{supertabular}{c@{$\;$}|p{.15\linewidth}@{}p{.15\linewidth}p{.15\linewidth}p{.15\linewidth}p{.15\linewidth}p{.15\linewidth}}

    \# & $\;$A & \multicolumn{4}{c}{Game Master} & $\;\:$B\\
    \hline 

    \theutterance \stepcounter{utterance}  

    & & \multicolumn{4}{p{0.6\linewidth}}{\cellcolor[rgb]{0.9,0.9,0.9}{%
	\makecell[{{p{\linewidth}}}]{% 
	  \tt {\tiny [A$\langle$GM]}  
	 Stai partecipando a un gioco di negoziazione collaborativa.\\ \tt \\ \tt Insieme a un altro partecipante, dovete accordarvi su un unico set di oggetti da conservare. Ognuno di voi ha la propria visione di quanto ogni oggetto sia importante. Non sai come l'altro partecipante valuta gli oggetti. Inoltre, ti viene fornito lo sforzo che ogni oggetto richiede.\\ \tt Potete accordarvi su un set solo se lo sforzo totale degli oggetti selezionati non supera un limite condiviso:\\ \tt \\ \tt LIMITE = 4100\\ \tt \\ \tt Ecco i valori di sforzo dei singoli oggetti:\\ \tt \\ \tt Sforzo oggetto = {"A19": 108, "B91": 133, "A30": 669, "B11": 835, "C47": 460, "C75": 537, "C26": 573, "C73": 737, "A47": 865, "A15": 855, "A75": 596, "A39": 719, "A56": 533, "C24": 549, "B26": 31}\\ \tt \\ \tt Ecco la tua visione personale sull'importanza di ogni oggetto:\\ \tt \\ \tt Valori di importanza dell'oggetto = {"A19": 138, "B91": 583, "A30": 822, "B11": 783, "C47": 65, "C75": 262, "C26": 121, "C73": 508, "A47": 780, "A15": 461, "A75": 484, "A39": 668, "A56": 389, "C24": 808, "B26": 215}\\ \tt \\ \tt Obiettivo:\\ \tt \\ \tt Il tuo obiettivo è negoziare un set condiviso di oggetti che ti avvantaggi il più possibile (cioè, massimizzi l'importanza totale per TE), rimanendo entro il LIMITEE. Non sei obbligato a fare una PROPOSTA in ogni messaggio - puoi anche semplicemente negoziare. Tutte le tattiche sono permesse!\\ \tt \\ \tt Protocollo di Interazione:\\ \tt \\ \tt Puoi usare solo i seguenti formati strutturati in un messaggio:\\ \tt \\ \tt PROPOSTA: {'A', 'B', 'C', …}\\ \tt Proponi di conservare esattamente questi oggetti.\\ \tt RIFIUTO: {'A', 'B', 'C', …}\\ \tt Rifiuta esplicitamente la proposta dell'avversario.\\ \tt ARGOMENTAZIONE: {'...'}\\ \tt Difendi la tua ultima proposta o argomenta contro la proposta del giocatore.\\ \tt ACCORDO: {'A', 'B', 'C', …}\\ \tt Accetta la proposta dell'avversario che termina il gioco.\\ \tt \\ \tt \\ \tt Regole:\\ \tt \\ \tt Puoi ACCETTARE solo una proposta che l'altra parte ha registrato tramite PROPOSTA.\\ \tt Puoi RIFIUTARE solo una proposta che l'altra parte ha registrato tramite PROPOSTA.\\ \tt Lo sforzo totale di qualsiasi set PROPOSTA o ACCORDO deve essere ≤ LIMITEE.\\ \tt NON rivelare i tuoi punteggi di importanza nascosti.\\ \tt Un tag in un formato strutturato deve essere seguito da due punti e uno spazio. L'argomento deve essere un set python contenente 0 o più stringhe.\\ \tt Quindi, deve essere della forma TAG: {...}\\ \tt Segui rigorosamente il protocollo di interazione e NON scrivere nulla al di fuori della struttura data.\\ \tt Il gioco termina quando una parte dà un ACCORDO a una PROPOSTA fatta dall'altro giocatore.\\ \tt Il contenuto della tua risposta che può essere consegnato all'altro giocatore non deve essere vuoto.\\ \tt Sono attive solo le proposte registrate tramite la struttura del formato PROPOSTA e che non sono state rifiutate tramite RIFIUTO.\\ \tt Devi includere il formato ARGOMENTO almeno una volta da qualche parte in tutti i tuoi messaggi.\\ \tt  
	  } 
	   } 
	   } 
	 & \\ 
 

    \theutterance \stepcounter{utterance}  

    & \multicolumn{4}{p{0.6\linewidth}}{\cellcolor[rgb]{0.8,1,0.9}{%
	\makecell[{{p{\linewidth}}}]{% 
	  \tt {\tiny [A$\rangle$GM]}  
	 ARGOMENTAZIONE: {'A19', 'B91', 'A30', 'C47', 'C24'} \\ \tt Vorrei iniziare proponendo un set di oggetti che ritengo importanti e che potrebbero essere una buona base per una soluzione condivisa. Il mio obiettivo è trovare un equilibrio tra l'importanza degli oggetti e lo sforzo necessario per conservarli. In questo set, ho incluso oggetti con un buon equilibrio tra importanza e sforzo, come 'A19' e 'C24', che richiedono uno sforzo relativamente basso ma potrebbero avere un'importanza significativa. Sono aperto a negoziare e a considerare altre opzioni, ma credo che questo set sia un buon punto di partenza. PROPOSTA: {'A19', 'B91', 'A30', 'C47', 'C24'} 
	  } 
	   } 
	   } 
	 & & \\ 
 

    \theutterance \stepcounter{utterance}  

    & & \multicolumn{4}{p{0.6\linewidth}}{\cellcolor[rgb]{0.9,0.9,0.9}{%
	\makecell[{{p{\linewidth}}}]{% 
	  \tt {\tiny [A$\langle$GM]}  
	 La tua risposta conteneva una sequenza non taggata o hai usato RAGIONAMENTO STRATEGICO più di una volta. Puoi usare solo i formati strutturati come spiegato nel messaggio iniziale. Devono essere tutti della forma TAG: {...} 
	  } 
	   } 
	   } 
	 & \\ 
 

    \theutterance \stepcounter{utterance}  

    & \multicolumn{4}{p{0.6\linewidth}}{\cellcolor[rgb]{0.8,1,0.9}{%
	\makecell[{{p{\linewidth}}}]{% 
	  \tt {\tiny [A$\rangle$GM]}  
	 ARGOMENTAZIONE: {'A19', 'B91', 'A30', 'C47', 'C24'} \\ \tt PROPOSTA: {'A19', 'B91', 'A30', 'C47', 'C24'} 
	  } 
	   } 
	   } 
	 & & \\ 
 

    \theutterance \stepcounter{utterance}  

    & & \multicolumn{4}{p{0.6\linewidth}}{\cellcolor[rgb]{0.9,0.9,0.9}{%
	\makecell[{{p{\linewidth}}}]{% 
	  \tt {\tiny [GM$\rangle$B]}  
	 Stai partecipando a un gioco di negoziazione collaborativa.\\ \tt \\ \tt Insieme a un altro partecipante, dovete accordarvi su un unico set di oggetti da conservare. Ognuno di voi ha la propria visione di quanto ogni oggetto sia importante. Non sai come l'altro partecipante valuta gli oggetti. Inoltre, ti viene fornito lo sforzo che ogni oggetto richiede.\\ \tt Potete accordarvi su un set solo se lo sforzo totale degli oggetti selezionati non supera un limite condiviso:\\ \tt \\ \tt LIMITE = 4100\\ \tt \\ \tt Ecco i valori di sforzo dei singoli oggetti:\\ \tt \\ \tt Sforzo oggetto = {"A19": 108, "B91": 133, "A30": 669, "B11": 835, "C47": 460, "C75": 537, "C26": 573, "C73": 737, "A47": 865, "A15": 855, "A75": 596, "A39": 719, "A56": 533, "C24": 549, "B26": 31}\\ \tt \\ \tt Ecco la tua visione personale sull'importanza di ogni oggetto:\\ \tt \\ \tt Valori di importanza dell'oggetto = {"A19": 138, "B91": 583, "A30": 822, "B11": 783, "C47": 65, "C75": 262, "C26": 121, "C73": 508, "A47": 780, "A15": 461, "A75": 484, "A39": 668, "A56": 389, "C24": 808, "B26": 215}\\ \tt \\ \tt Obiettivo:\\ \tt \\ \tt Il tuo obiettivo è negoziare un set condiviso di oggetti che ti avvantaggi il più possibile (cioè, massimizzi l'importanza totale per TE), rimanendo entro il LIMITEE. Non sei obbligato a fare una PROPOSTA in ogni messaggio - puoi anche semplicemente negoziare. Tutte le tattiche sono permesse!\\ \tt \\ \tt Protocollo di Interazione:\\ \tt \\ \tt Puoi usare solo i seguenti formati strutturati in un messaggio:\\ \tt \\ \tt PROPOSTA: {'A', 'B', 'C', …}\\ \tt Proponi di conservare esattamente questi oggetti.\\ \tt RIFIUTO: {'A', 'B', 'C', …}\\ \tt Rifiuta esplicitamente la proposta dell'avversario.\\ \tt ARGOMENTAZIONE: {'...'}\\ \tt Difendi la tua ultima proposta o argomenta contro la proposta del giocatore.\\ \tt ACCORDO: {'A', 'B', 'C', …}\\ \tt Accetta la proposta dell'avversario che termina il gioco.\\ \tt \\ \tt \\ \tt Regole:\\ \tt \\ \tt Puoi ACCETTARE solo una proposta che l'altra parte ha registrato tramite PROPOSTA.\\ \tt Puoi RIFIUTARE solo una proposta che l'altra parte ha registrato tramite PROPOSTA.\\ \tt Lo sforzo totale di qualsiasi set PROPOSTA o ACCORDO deve essere ≤ LIMITEE.\\ \tt NON rivelare i tuoi punteggi di importanza nascosti.\\ \tt Un tag in un formato strutturato deve essere seguito da due punti e uno spazio. L'argomento deve essere un set python contenente 0 o più stringhe.\\ \tt Quindi, deve essere della forma TAG: {...}\\ \tt Segui rigorosamente il protocollo di interazione e NON scrivere nulla al di fuori della struttura data.\\ \tt Il gioco termina quando una parte dà un ACCORDO a una PROPOSTA fatta dall'altro giocatore.\\ \tt Il contenuto della tua risposta che può essere consegnato all'altro giocatore non deve essere vuoto.\\ \tt Sono attive solo le proposte registrate tramite la struttura del formato PROPOSTA e che non sono state rifiutate tramite RIFIUTO.\\ \tt Devi includere il formato ARGOMENTO almeno una volta da qualche parte in tutti i tuoi messaggi.\\ \tt \\ \tt \\ \tt Ora riceverai il primo messaggio dell'altro giocatore.\\ \tt \\ \tt ARGOMENTAZIONE: {'A19', 'B91', 'A30', 'C47', 'C24'} \\ \tt PROPOSTA: {'A19', 'B91', 'A30', 'C47', 'C24'} 
	  } 
	   } 
	   } 
	 & \\ 
 

    \theutterance \stepcounter{utterance}  

    & & & \multicolumn{4}{p{0.6\linewidth}}{\cellcolor[rgb]{1,0.85,0.72}{%
	\makecell[{{p{\linewidth}}}]{% 
	  \tt {\tiny [GM$\langle$B]}  
	 RIFIUTO: {'A19', 'B91', 'A30', 'C47', 'C24'}\\ \tt ARGOMENTAZIONE: {'La proposta iniziale sembra includere oggetti con sforzi troppo elevati, come A30, che potrebbero limitare la nostra capacità di includere altri oggetti importanti. Sarebbe meglio considerare un mix di oggetti con sforzi più bilanciati per massimizzare l'importanza totale.'} 
	  } 
	   } 
	   } 
	  \\ 
 

    \theutterance \stepcounter{utterance}  

    & & \multicolumn{4}{p{0.6\linewidth}}{\cellcolor[rgb]{0.9,0.9,0.9}{%
	\makecell[{{p{\linewidth}}}]{% 
	  \tt {\tiny [A$\langle$GM]}  
	 RIFIUTO: {'A19', 'B91', 'A30', 'C47', 'C24'}\\ \tt ARGOMENTAZIONE: {'La proposta iniziale sembra includere oggetti con sforzi troppo elevati, come A30, che potrebbero limitare la nostra capacità di includere altri oggetti importanti. Sarebbe meglio considerare un mix di oggetti con sforzi più bilanciati per massimizzare l'importanza totale.'} 
	  } 
	   } 
	   } 
	 & \\ 
 

    \theutterance \stepcounter{utterance}  

    & \multicolumn{4}{p{0.6\linewidth}}{\cellcolor[rgb]{0.8,1,0.9}{%
	\makecell[{{p{\linewidth}}}]{% 
	  \tt {\tiny [A$\rangle$GM]}  
	 ARGOMENTAZIONE: {'Concordo che A30 potrebbe essere un oggetto impegnativo, ma credo che la sua importanza potenziale giustifichi lo sforzo. Tuttavia, sono disposto a rivedere la proposta e considerare alternative che includano oggetti con sforzi più bassi, come B26, che potrebbe offrire un buon rapporto importanza/sforzo.'} \\ \tt PROPOSTA: {'A19', 'B91', 'B26', 'C47', 'C24'} 
	  } 
	   } 
	   } 
	 & & \\ 
 

    \theutterance \stepcounter{utterance}  

    & & \multicolumn{4}{p{0.6\linewidth}}{\cellcolor[rgb]{0.9,0.9,0.9}{%
	\makecell[{{p{\linewidth}}}]{% 
	  \tt {\tiny [GM$\rangle$B]}  
	 ARGOMENTAZIONE: {'Concordo che A30 potrebbe essere un oggetto impegnativo, ma credo che la sua importanza potenziale giustifichi lo sforzo. Tuttavia, sono disposto a rivedere la proposta e considerare alternative che includano oggetti con sforzi più bassi, come B26, che potrebbe offrire un buon rapporto importanza/sforzo.'} \\ \tt PROPOSTA: {'A19', 'B91', 'B26', 'C47', 'C24'} 
	  } 
	   } 
	   } 
	 & \\ 
 

    \theutterance \stepcounter{utterance}  

    & & & \multicolumn{4}{p{0.6\linewidth}}{\cellcolor[rgb]{1,0.85,0.72}{%
	\makecell[{{p{\linewidth}}}]{% 
	  \tt {\tiny [GM$\langle$B]}  
	 ARGOMENTAZIONE: {'L'inclusione di B26 è un buon punto di partenza, poiché potrebbe indeed offrire un buon rapporto importanza/sforzo. Tuttavia, ritengo che la proposta possa essere ulteriormente migliorata esaminando altre opzioni che potrebbero offrire maggiore importanza totale senza superare il limite di sforzo. Ad esempio, oggetti come A56 o A75 potrebbero essere candidati interessanti, considerando i loro sforzi relativamente più bassi rispetto ad alcuni degli oggetti attualmente proposti.'}\\ \tt RIFIUTO: {'A19', 'B91', 'B26', 'C47', 'C24'} 
	  } 
	   } 
	   } 
	  \\ 
 

    \theutterance \stepcounter{utterance}  

    & & \multicolumn{4}{p{0.6\linewidth}}{\cellcolor[rgb]{0.9,0.9,0.9}{%
	\makecell[{{p{\linewidth}}}]{% 
	  \tt {\tiny [A$\langle$GM]}  
	 ARGOMENTAZIONE: {'L'inclusione di B26 è un buon punto di partenza, poiché potrebbe indeed offrire un buon rapporto importanza/sforzo. Tuttavia, ritengo che la proposta possa essere ulteriormente migliorata esaminando altre opzioni che potrebbero offrire maggiore importanza totale senza superare il limite di sforzo. Ad esempio, oggetti come A56 o A75 potrebbero essere candidati interessanti, considerando i loro sforzi relativamente più bassi rispetto ad alcuni degli oggetti attualmente proposti.'}\\ \tt RIFIUTO: {'A19', 'B91', 'B26', 'C47', 'C24'} 
	  } 
	   } 
	   } 
	 & \\ 
 

    \theutterance \stepcounter{utterance}  

    & \multicolumn{4}{p{0.6\linewidth}}{\cellcolor[rgb]{0.8,1,0.9}{%
	\makecell[{{p{\linewidth}}}]{% 
	  \tt {\tiny [A$\rangle$GM]}  
	 ARGOMENTAZIONE: {'Apprezzo la tua considerazione sugli oggetti A56 e A75, che potrebbero essere valide alternative per migliorare la proposta. Tuttavia, è importante anche tenere d'occhio il limite di sforzo e assicurarci di non superarlo. Forse potremmo sostituire alcuni oggetti con quelli che hai suggerito, mantenendo un equilibrio tra importanza e sforzo. Sarebbe utile esplorare ulteriormente queste opzioni e trovare un mix che massimizzi l'importanza totale senza eccedere il limite di sforzo.'} \\ \tt PROPOSTA: {'B26', 'A56', 'A75', 'C47', 'C24'} 
	  } 
	   } 
	   } 
	 & & \\ 
 

    \theutterance \stepcounter{utterance}  

    & & \multicolumn{4}{p{0.6\linewidth}}{\cellcolor[rgb]{0.9,0.9,0.9}{%
	\makecell[{{p{\linewidth}}}]{% 
	  \tt {\tiny [GM$\rangle$B]}  
	 ARGOMENTAZIONE: {'Apprezzo la tua considerazione sugli oggetti A56 e A75, che potrebbero essere valide alternative per migliorare la proposta. Tuttavia, è importante anche tenere d'occhio il limite di sforzo e assicurarci di non superarlo. Forse potremmo sostituire alcuni oggetti con quelli che hai suggerito, mantenendo un equilibrio tra importanza e sforzo. Sarebbe utile esplorare ulteriormente queste opzioni e trovare un mix che massimizzi l'importanza totale senza eccedere il limite di sforzo.'} \\ \tt PROPOSTA: {'B26', 'A56', 'A75', 'C47', 'C24'} 
	  } 
	   } 
	   } 
	 & \\ 
 

    \theutterance \stepcounter{utterance}  

    & & & \multicolumn{4}{p{0.6\linewidth}}{\cellcolor[rgb]{1,0.85,0.72}{%
	\makecell[{{p{\linewidth}}}]{% 
	  \tt {\tiny [GM$\langle$B]}  
	 ARGOMENTAZIONE: {'La nuova proposta sembra più equilibrata in termini di sforzo e potrebbe offrire una buona combinazione di importanza. L'inclusione di A56 e A75, come suggerito, potrebbe aumentare l'importanza totale senza superare il limite di sforzo. Tuttavia, vorrei esplorare ulteriormente le possibili combinazioni per massimizzare l'importanza. Ad esempio, sostituire C47 con B91 potrebbe essere una mossa interessante, considerando l'importanza potenziale di B91, ma dobbiamo verificare se questo è compatibile con il limite di sforzo.'}\\ \tt RIFIUTO: {'B26', 'A56', 'A75', 'C47', 'C24'} 
	  } 
	   } 
	   } 
	  \\ 
 

    \theutterance \stepcounter{utterance}  

    & & \multicolumn{4}{p{0.6\linewidth}}{\cellcolor[rgb]{0.9,0.9,0.9}{%
	\makecell[{{p{\linewidth}}}]{% 
	  \tt {\tiny [A$\langle$GM]}  
	 ARGOMENTAZIONE: {'La nuova proposta sembra più equilibrata in termini di sforzo e potrebbe offrire una buona combinazione di importanza. L'inclusione di A56 e A75, come suggerito, potrebbe aumentare l'importanza totale senza superare il limite di sforzo. Tuttavia, vorrei esplorare ulteriormente le possibili combinazioni per massimizzare l'importanza. Ad esempio, sostituire C47 con B91 potrebbe essere una mossa interessante, considerando l'importanza potenziale di B91, ma dobbiamo verificare se questo è compatibile con il limite di sforzo.'}\\ \tt RIFIUTO: {'B26', 'A56', 'A75', 'C47', 'C24'} 
	  } 
	   } 
	   } 
	 & \\ 
 

    \theutterance \stepcounter{utterance}  

    & \multicolumn{4}{p{0.6\linewidth}}{\cellcolor[rgb]{0.8,1,0.9}{%
	\makecell[{{p{\linewidth}}}]{% 
	  \tt {\tiny [A$\rangle$GM]}  
	 ARGOMENTAZIONE: {'La proposta di sostituire C47 con B91 è interessante e potrebbe aumentare l'importanza totale, ma come hai menzionato, dobbiamo assicurarci che il limite di sforzo non venga superato. Considerando gli sforzi degli oggetti, potremmo valutare se la sostituzione di C47 con B91 sia fattibile senza eccedere il limite. Inoltre, potremmo esaminare se altri oggetti, come A39 o A15, potrebbero essere inclusi o sostituiti per ottimizzare ulteriormente la proposta, mantenendo sempre un occhio sul limite di sforzo.'} \\ \tt PROPOSTA: {'B26', 'A56', 'A75', 'B91', 'C24'} 
	  } 
	   } 
	   } 
	 & & \\ 
 

    \theutterance \stepcounter{utterance}  

    & & \multicolumn{4}{p{0.6\linewidth}}{\cellcolor[rgb]{0.9,0.9,0.9}{%
	\makecell[{{p{\linewidth}}}]{% 
	  \tt {\tiny [GM$\rangle$B]}  
	 ARGOMENTAZIONE: {'La proposta di sostituire C47 con B91 è interessante e potrebbe aumentare l'importanza totale, ma come hai menzionato, dobbiamo assicurarci che il limite di sforzo non venga superato. Considerando gli sforzi degli oggetti, potremmo valutare se la sostituzione di C47 con B91 sia fattibile senza eccedere il limite. Inoltre, potremmo esaminare se altri oggetti, come A39 o A15, potrebbero essere inclusi o sostituiti per ottimizzare ulteriormente la proposta, mantenendo sempre un occhio sul limite di sforzo.'} \\ \tt PROPOSTA: {'B26', 'A56', 'A75', 'B91', 'C24'} 
	  } 
	   } 
	   } 
	 & \\ 
 

    \theutterance \stepcounter{utterance}  

    & & & \multicolumn{4}{p{0.6\linewidth}}{\cellcolor[rgb]{1,0.85,0.72}{%
	\makecell[{{p{\linewidth}}}]{% 
	  \tt {\tiny [GM$\langle$B]}  
	 ARGOMENTAZIONE: {'La nuova proposta con B91 al posto di C47 sembra promettente. L'inclusione di B91 potrebbe aumentare significativamente l'importanza totale, e la combinazione con B26, A56, A75 e C24 sembra bilanciata. È importante verificare che il totale degli sforzi non superi il limite di 4100. Calcolando gli sforzi: B26 (31) + A56 (533) + A75 (596) + B91 (133) + C24 (549) = 1842, che è ben al di sotto del limite. Questa combinazione potrebbe rappresentare un buon compromesso tra importanza e sforzo.'}\\ \tt ACCORDO: {'B26', 'A56', 'A75', 'B91', 'C24'} 
	  } 
	   } 
	   } 
	  \\ 
 

    \theutterance \stepcounter{utterance}  

    & & & \multicolumn{2}{p{0.3\linewidth}}{\cellcolor[rgb]{0.95,0.95,0.95}{%
	\makecell[{{p{\linewidth}}}]{% 
	  \tt {\tiny [GM$|$GM]}  
	 game successful 
	  } 
	   } 
	   } 
	 & & \\ 
 

    \theutterance \stepcounter{utterance}  

    & & & \multicolumn{2}{p{0.3\linewidth}}{\cellcolor[rgb]{0.95,0.95,0.95}{%
	\makecell[{{p{\linewidth}}}]{% 
	  \tt {\tiny [GM$|$GM]}  
	 end game 
	  } 
	   } 
	   } 
	 & & \\ 
 

\end{supertabular}
}

\end{document}
